\subsection{Germs of holomorphic functions}

Let E and F be complex Banach spaces. We recall (cf. VAR, R1, p.~26, 3.2.1, p.~22, 3.1 and p.~88, App.) that a holomorphic mapping defined on an open subset U of E and taking values in F is a mapping \(f \colon \mathrm{U} \to \mathrm{F}\) such that, for all \(x \in \mathrm{U}\) there exists a convergent series

\[
f _ {x} = \sum_ {k \geqslant 0} f _ {x, k}
\]

satisfying \(f(x + y) = f_x(y)\) for all \(y \in E\) sufficiently close to 0, where \(f_{x,k} \colon E \to C\) is a continuous homogeneous polynomial of degree k on E with values in F, that is, a mapping of the form

\[
f _ {x, k} (y) = \widetilde {f} _ {x, k} (y, \ldots , y)
\]

where \(\widetilde{f}_{x,k}\colon\mathrm{E}^{k}\to\mathrm{F}\) is a map \(k\) continuous multilinear. We denote \(\mathcal{O}(\mathrm{U};\mathrm{F})\) the complex vector space of holomorphic functions on U with values in F, equipped with the topology of compact convergence. It is a locally convex topological vector space, whose topology is defined by the seminorms \(f\mapsto\sup_{z\in\mathrm{K}}\|f(z)\|\) where K ranges over the set of compact subsets of U.

Let G be a complex Banach space and V an open subset of G. For every holomorphic map \(\varphi\colon\mathrm{V}\to\mathrm{U}\) , the map \(\varphi^{*}\colon f\mapsto f\circ\varphi\) is a continuous linear map from \(\mathcal{O}(\mathrm{U};\mathrm{F})\) into \(\mathcal{O}(\mathrm{V};\mathrm{F})\) .

If H is a complex Banach space and \(\varphi\colon\mathrm{F}\to\mathrm{H}\) a continuous linear map, then the map \(f\mapsto\varphi\circ f\) is a continuous linear map from \(\mathcal{O}(\mathrm{U};\mathrm{F})\) into \(\mathcal{O}(\mathrm{U};\mathrm{H})\) , denoted \(\varphi_{*}\) .

Let n be a natural number and set \(E = C^{n}\) . Let K be a compact subset of \(C^{n}\) and let U be the decreasing filtered set of open neighborhoods of K. If U, \(U' \in U\) and \(U' \subset U\) , the restriction map of functions from \(\mathcal{O}(U; F)\) into \(\mathcal{O}(U'; F)\) is continuous. The inductive limit of the spaces \(\mathcal{O}(\mathrm{U};\mathrm{F})\) for these maps is denoted \(\mathcal{O}(\mathrm{K};\mathrm{F})\) . The elements of \(\mathcal{O}(\mathrm{K};\mathrm{F})\) are called the germs of holomorphic functions in a neighborhood of K with values in F.

The space \(\mathcal{O}(\mathrm{K};\mathrm{F})\) is endowed with the inductive limit topology of the locally convex topologies of the \(\mathcal{O}(\mathrm{U};\mathrm{F})\) Let X be a locally convex topological vector space and \(\varphi\colon\mathcal{O}(\mathrm{K};\mathrm{F})\to\mathrm{X}\) a map. For every open neighborhood U of K, the map \(\mathcal{O}(\mathrm{U};\mathrm{F})\to\mathrm{X}\) derived from \(\varphi\) by composition with the canonical map \(\mathcal{O}(\mathrm{U};\mathrm{F})\to\mathcal{O}(\mathrm{K};\mathrm{F})\) is denoted \(\varphi^{\mathrm{U}}\) . The map \(\varphi\) is continuous if and only if \(\varphi^{\mathrm{U}}\) is continuous for every U (EVT, II, p.~29, prop. 5, (iii)).

Let m be a natural number. Let L be a compact subset of \(C^{m}\) and let V be an open neighborhood of L. Let \(\varphi\colon V \to C^{n}\) a holomorphic map such that \(\varphi(L) \subset K\) . The continuous linear maps

\[
\mathcal {O} (\mathrm{U}; \mathrm{F}) \xrightarrow {\varphi^ {*}} \mathcal {O} (\stackrel {- 1} {\varphi} (\mathrm{U}); \mathrm{F}) \xrightarrow {\varphi^ {- 1} \stackrel {- 1} {\varphi} (\mathrm{U})} \mathcal {O} (\mathrm{L}; \mathrm{F})
\]

for U an open neighborhood of K, induce a continuous linear map \(\varphi^{*}\colon\mathcal{O}(\mathrm{K};\mathrm{F})\to\mathcal{O}(\mathrm{L};\mathrm{F})\) (loc. cit.).

Let H be a complex Banach space and \(\varphi\colon\mathrm{F}\to\mathrm{H}\) a continuous linear map. The continuous linear maps

\[
\mathcal {O} (\mathrm{U}; \mathrm{F}) \xrightarrow {\varphi_ {*}} \mathcal {O} (\mathrm{U}; \mathrm{H}) \xrightarrow {\varphi^ {\mathrm{U}}} \mathcal {O} (\mathrm{K}; \mathrm{F})
\]

where U ranges over the set of open neighborhoods of K in \(C^{n}\) induce a continuous linear map \(\varphi_{*}\) of \(\mathcal{O}(K;F)\) into \(\mathcal{O}(K;H)\) (loc. cit.). We shall sometimes denote \(\varphi \circ f = \varphi_{*}(f)\) .

For every open neighborhood U of K, the restriction to K is a continuous linear map. \(\mathcal{O}(\mathrm{U};\mathrm{F})\to \mathcal{C}(\mathrm{K};\mathrm{F})\) ; these maps induce a continuous linear map \(\mathcal{O}(\mathrm{K};\mathrm{F})\to \mathcal{C}(\mathrm{K};\mathrm{F})\) called evaluation of germs of holomorphic functions on K.

Let A be a unital complex Banach algebra. The spaces \(\mathcal{O}(\mathrm{U};\mathrm{A})\) and \(\mathcal{O}(\mathrm{K};\mathrm{A})\) are unital algebras. If \(A \neq \{0\}\) one can canonically identify \(\mathcal{O}(\mathrm{U};\mathbf{C})\) (resp. \(\mathcal{O}(\mathrm{K};\mathbf{C})\) ) with the subalgebra \(\mathcal{O}(\mathrm{U};\mathbf{C}) \cdot 1\) of \(\mathcal{O}(\mathrm{U};\mathrm{A})\) (resp. to the subalgebra \(\mathcal{O}(\mathrm{K};\mathbf{C}) \cdot 1\) of \(\mathcal{O}(\mathrm{K};\mathrm{A})\) ). We shall set \(\mathcal{O}(\mathrm{U}) = \mathcal{O}(\mathrm{U};\mathbf{C})\) and \(\mathcal{O}(\mathrm{K}) = \mathcal{O}(\mathrm{K};\mathbf{C})\) .

